\chapter{Quantenbits, Schrödinger-Gleichung und Dirac-Gleichung in der fraktalen T0-Geometrie}


	
	\subsection*{Narrative Einführung: Das kosmische Gehirn im Detail}
	
	Wir setzen unsere Reise durch das kosmische Gehirn fort. In diesem Kapitel betrachten wir weitere Aspekte der fraktalen Struktur des Universums, die – wie die komplexen Windungen eines Gehirns – auf allen Skalen selbstähnliche Muster aufweisen. Was auf den ersten Blick wie isolierte physikalische Phänomene erscheint, erweist sich bei genauerer Betrachtung als Ausdruck eines einheitlichen geometrischen Prinzips. Hier leiten wir Quantenbits, die Schrödinger-Gleichung und eine alternative konforme Beschreibung der Dirac-Gleichung aus der Dynamik des Vakuumphasenfeldes ab – ohne 4×4-Matrixstruktur, da Spin topologisch aus dem Vakuumfeld emergiert.
	
	\subsection*{Mathematische Grundlage}
	
	Quantencomputing und relativistische Quantenmechanik basieren auf Zustandsüberlagerung und Spin. In der FFGFT (fraktalen T0-Geometrie) emergieren diese aus Phasenmoden des Vakuumfeldes \(\theta(x,t)\), reguliert durch \(\xipar = \frac{4}{3} \times 10^{-4}\). Die Wellenfunktion \(\psi\) ist ein rein mathematisches Hilfskonstrukt – keine ontologische Realität. Die Dirac-Gleichung erhält in T0 eine alternative konforme Beschreibung als skalare Feldknoten-Dynamik, wobei die 4×4-Matrixstruktur durch topologische Phasenwindungen ersetzt wird. Alle experimentellen Vorhersagen bleiben erhalten.
	
	
	
	\subsection*{Qubit als lokale Phasenexcitation}
	
	Ein Quantenbit (Qubit) wird durch eine kleine lokale Abweichung der Vakuumphase dargestellt:
	
	\begin{equation}
		\delta \theta = \alpha \cdot 0 + \beta \cdot \pi.
	\end{equation}
	
	Diese Gleichung definiert die beiden Basiszustände: \(\ket{0}\) entspricht keiner Phasenabweichung (\(\delta \theta = 0\)), \(\ket{1}\) einer Phasenverschiebung um \(\pi\). Die Koeffizienten \(\alpha\) und \(\beta\) sind komplexe Amplituden mit \(|\alpha|^2 + |\beta|^2 = 1\), die die Wahrscheinlichkeiten für Messungen kodieren. Die Superposition:
	
	\begin{equation}
		\ket{\psi} = \alpha \ket{0} + \beta \ket{1}
	\end{equation}
	
	ist ein rein mathematisches Hilfskonstrukt zur Beschreibung möglicher Messergebnisse. Ontologisch existiert nur die deterministische Phase \(\theta + \delta \theta\) des Vakuumfeldes – keine reale Überlagerung von Zuständen.
	
	\textbf{Einheitenprüfung:}
	\[
	\begin{aligned}
		[\delta \theta] &= \text{dimensionslos}.
	\end{aligned}
	\]
	
	\subsection*{Emergenz der Schrödinger-Gleichung}
	
	Die nicht-relativistische Dynamik eines Teilchens folgt aus der Phasenentwicklung. Die fraktale T0-Geometrie erweitert die übliche Schrödinger-Gleichung um einen Phasenkohärenz-Term:
	
	\begin{equation}
		i \hbar \partial_t \psi = \left( -\frac{\hbar^2}{2m} \nabla^2 + V(x) \right) \psi \cdot (1 + \xipar \ln(l/l_0)).
	\end{equation}
	
	Die linke Seite beschreibt die zeitliche Entwicklung der Wellenfunktion \(\psi\) (emergent aus der Phase \(\theta\)). Der kinetische Term \(-\frac{\hbar^2}{2m} \nabla^2\) entsteht aus der lokalen Krümmung der Phase, das Potenzial \(V(x)\) aus lokaler Amplitude-Deformation. Der Faktor \((1 + \xipar \ln(l/l_0))\) ist die fraktale Erweiterung: Er berücksichtigt die logarithmische Phasenkohärenz über Skalen \(l\) relativ zur Korrelationslänge \(l_0\). Für \(l \ll l_0 \xipar^{-1}\) reduziert sich die Gleichung exakt auf die Standard-Schrödinger-Gleichung – die FFGFT ist eine natürliche Erweiterung der üblichen Form.
	
	\textbf{Einheitenprüfung:}
	\[
	\begin{aligned}
		[i \hbar \partial_t \psi] &= \si{\joule\second} \cdot \si{\per\second} = \si{\joule}.
	\end{aligned}
	\]
	
	\subsection*{Emergenz der Dirac-Gleichung}
	
	Die Standard-Dirac-Gleichung \(i \hbar \gamma^\mu \partial_\mu \psi - m c \psi = 0\) emergiert in T0 aus multi-komponentigen Vakuumfeldern. In natürlichen Einheiten (\(\hbar = c = 1\)) lautet die modifizierte Dirac-Gleichung:
	
	\begin{equation}
		i\gamma^{\mu}(\partial_{\mu} + \Gamma_{\mu}^{(T)}) \psi - m(\vec{x},t) \psi = 0,
	\end{equation}
	
	wobei \(\Gamma_{\mu}^{(T)} = -\partial_{\mu} m / m^2\) die Time-Field-Verbindung ist, die aus der Time-Mass-Dualität \(T \cdot m = 1\) folgt, und \(m(\vec{x},t)\) die lokale Massendichte. Im schwachen Feld-Limit (\(\beta \ll 1\)) reduziert sich die Gleichung exakt auf die Standard-Dirac-Gleichung, konsistent mit QED-Präzisionsmessungen.
	
	\subsection*{Alternative konforme Beschreibung: Dirac-Gleichung als Feldknoten-Dynamik}
	
	In der T0-Perspektive wird die Dirac-Gleichung konform als skalare Wellengleichung für Massefeldknoten beschrieben. Dies ist keine Vereinfachung, sondern eine alternative Formulierung, die ohne die 4×4-Matrixstruktur auskommt, da die Spin-Information topologisch im Vakuumfeld kodiert ist:
	
	\begin{equation}
		\square \, \delta m = 0,
	\end{equation}
	
	wobei \(\square\) der D'Alembert-Operator und \(\delta m\) die Feldknoten-Amplitude (Abweichung vom Vakuumgrund \(\rho_0\)) ist. Der Spinor \(\psi\) wird zu einem Knotenmuster \(\psi(x,t) \to \delta m_0 \cdot f_{\text{spin}}(x,t)\), wobei die Spin-Strukturfunktion \(f_{\text{spin}}\) die halbzahlige Phasenwindung kodiert. Spin-1/2 emergiert aus der Knotenrotation mit Frequenz \(\omega_{\text{spin}} \propto m c^2 / \hbar \cdot \xipar\). Die zugehörige Lagrangedichte lautet \(\mathcal{L} = \varepsilon \cdot (\partial \, \delta m)^2\).
	
	\subsection*{Qubit-Gatter als Phasenmanipulationen}
	
	Beispiel Hadamard-Gate:
	
	\begin{equation}
		H: \delta \theta \to \frac{\delta \theta + \pi/2}{\sqrt{2}}.
	\end{equation}
	
	Es rotiert die Phase um 90° und normiert – erzeugt gleichmäßige Superposition. Alle Gatter sind lokale Phasenoperationen am Vakuumfeld.
	
	\subsection*{Vergleich Standard – FFGFT}
	
	\begin{center}
		\small
		\resizebox{\textwidth}{!}{%
			\begin{tabular}{p{0.45\textwidth}p{0.45\textwidth}}
				\textbf{Standard} & \textbf{FFGFT (T0)} \\
				\hline
				Qubit postuliert & Lokale Phasenexcitation \\
				Schrödinger/Dirac axiomatisch & Erweiterte Form / alternative konforme Beschreibung \\
				Spin ad-hoc & Topologische Windung der Phase \\
				Keine Gravitation & Einheitlich mit Amplitude \\
				Ontologische Superposition & Mathematisches Hilfskonstrukt \\
			\end{tabular}%
		}
	\end{center}
	
	\subsection*{Schlussfolgerung}
	
	Die FFGFT leitet Quantenbits als Phasenexcitationen, die Schrödinger-Gleichung als erweiterte nicht-relativistische Form und die Dirac-Gleichung als alternative konforme Beschreibung ohne 4×4-Matrizen ab. Superposition und Wellenfunktion sind rein mathematische Hilfskonstrukte – das Vakuumfeld bleibt deterministisch. Spin ist topologische Eigenschaft der Phase. Alles emergiert parameterfrei aus \(\xipar\), vereinheitlicht Quantencomputing mit fundamentaler Physik.



